\documentclass[11pt]{article}

\usepackage[margin=1in]{geometry}
\usepackage{amsmath,amssymb,amsthm,mathtools}
\usepackage{booktabs,array}
\usepackage{microtype}
\usepackage[numbers,sort&compress]{natbib}
\usepackage[hidelinks]{hyperref}
\usepackage[nameinlink,capitalise]{cleveref}

\newtheorem{theorem}{Theorem}[section]
\newtheorem{lemma}[theorem]{Lemma}
\newtheorem{corollary}[theorem]{Corollary}
\newtheorem{proposition}[theorem]{Proposition}
\theoremstyle{definition}
\newtheorem{remark}[theorem]{Remark}

\newcommand{\R}{\mathbb{R}}
\newcommand{\eps}{\varepsilon}
\newcommand{\one}{\mathbf{1}}

\title{Consequences of Rybin's Counterexample:\\
Convex-Decomposition Failure and a Planar $9/8$ Lower Bound\\
for Cost-Preserving Unsplittable Flow}
\author{Angus Muffatti}
\date{24 July 2026}

\begin{document}
\maketitle

\begin{abstract}
Rybin recently disproved Goemans' cost-preserving strengthening of the
Dinitz--Garg--Goemans theorem with a seven-vertex planar acyclic instance
whose fractional flow costs $58$ while every unsplittable flow within the
additive max-demand bound costs at least $60$ \citep{Rybin2026}.  We record
two consequences.  First, the same instance refutes the
cost-and-two-sided-bounds conjecture and the convex-decomposition
strengthening of Morell and Skutella; their costless two-sided existence
conjecture is untouched.  Second, a parametric form of the instance shows
that every universal cost-preserving rounding theorem on planar acyclic
digraphs requires additive congestion at least $(9/8 - o(1))\,D$, where $D$
is the maximum demand.  Combined with the planar factor-two theorem of
Traub, Vargas Koch, and Zenklusen, the optimal planar cost-preserving
constant lies in $[9/8,\,2]$.  The underlying mechanism is a violated
triangle stable-set facet, which we state as a reusable design principle.
\end{abstract}

\section{Introduction}

The theorem of Dinitz, Garg, and Goemans states that a feasible fractional
single-source flow can always be rounded to an unsplittable flow while
increasing every arc load by at most the maximum demand
\cite{DinitzGargGoemans1999}.  Goemans proposed that the same additive bound
should remain achievable with no increase in cost; the conjecture was
recorded in the minimum-cost literature by Skutella \cite{Skutella2002} and
guided work on min-cost unsplittable flow for roughly a quarter-century
\cite{MartensSalazarSkutella2007,MorellSkutella2022,SwamyTraubVargasKochZenklusen2026}.

In July 2026 Dmitry Rybin disproved the conjecture with an explicit
seven-vertex planar instance, constructed in a public AI-assisted research
session \citep{Rybin2026}.  \Cref{sec:instance} reproduces that instance and
its certificate, with full credit to Rybin, because the two contributions of
this note are read off directly from it: the corollaries in
\Cref{sec:corollaries} and the quantitative planar lower bound in
\Cref{sec:family}.

\section{Preliminaries}

Let $G=(V,A)$ be a directed acyclic graph with common source $s$, terminals
$t_1,\dots,t_k$, and demands $d_1,\dots,d_k>0$.  A fractional flow
$x\in\R_{\ge0}^{A}$ satisfies the corresponding divergence constraints.  An
unsplittable flow chooses one directed $s$--$t_i$ path $P_i$ per terminal and
has load
\[
 y_a=\sum_{i:a\in P_i} d_i \qquad (a\in A).
\]
Write $D:=\max_i d_i$.  The Dinitz--Garg--Goemans theorem guarantees an
unsplittable $y$ with $y_a\le x_a+D$ for all arcs
\cite{DinitzGargGoemans1999}.  Goemans' cost conjecture asked, for every
nonnegative arc-cost vector $c$, for such a $y$ satisfying additionally
$c^\top y\le c^\top x$.

\section{Rybin's instance}\label{sec:instance}

Everything in this section is due to Rybin \citep{Rybin2026}; we reproduce it
to keep the note self-contained.  Let
\[
 V=\{s,u,v,w,t_1,t_2,t_3\},
 \qquad
 (d_1,d_2,d_3)=(15,10,15),
 \qquad D=15,
\]
with arcs, fractional loads, and costs as in \Cref{tab:arcs}.

\begin{table}[ht]
\centering
\caption{Rybin's flow instance \citep{Rybin2026}.}
\label{tab:arcs}
\begin{tabular}{@{}lrr@{}}
\toprule
arc $a$ & $x_a$ & $c_a$\\
\midrule
$s\to t_1$ & 10 & 2\\
$s\to t_2$ & 6 & 3\\
$s\to u$   & 24 & 0\\
$u\to t_3$ & 10 & 2\\
$u\to v$   & 14 & 0\\
$v\to t_1$ & 5 & 0\\
$v\to w$   & 9 & 0\\
$w\to t_2$ & 4 & 0\\
$w\to t_3$ & 5 & 0\\
\bottomrule
\end{tabular}
\end{table}

Each terminal has exactly two source--terminal paths: the expensive direct
options $E_1=st_1$, $E_2=st_2$, $E_3=sut_3$ (each of total integral cost $30$
after multiplication by its demand) and the zero-cost options $Z_1=suvt_1$,
$Z_2=suvwt_2$, $Z_3=suvwt_3$.  The fractional flow is feasible with cost
$c^\top x=58$, and a routing satisfies $y\le x+D\one$ if and only if it uses
at most one of $Z_1,Z_2,Z_3$: each pair of zero-cost choices overloads one of
the arcs $su$, $uv$, $vw$ by exactly one unit beyond $x_a+D$.  Hence every
capacity-good routing uses at least two expensive paths and costs at least
$60$, and $(E_1,E_2,Z_3)$ attains $60$:
\[
 \min\{c^\top y:\ y\ \text{unsplittable},\ y\le x+D\one\}=60>58=c^\top x .
\]

\section{Two conjectures of Morell and Skutella fall with it}\label{sec:corollaries}

Morell and Skutella conjectured that every fractional single-source flow is a
convex combination of unsplittable flows satisfying the simultaneous bounds
$x_a-D\le y_a\le x_a+D$ on every arc, and noted that such a decomposition
would provide no-more-expensive members for arbitrary arc costs
\cite[Conjecture~2]{MorellSkutella2022}.

\begin{corollary}\label{cor:convex}
Rybin's instance refutes both the cost-and-two-sided-bounds conjecture and
the convex-decomposition strengthening of Morell and Skutella.  It does not
refute their costless two-sided existence conjecture.
\end{corollary}
\begin{proof}
Every two-sided-good routing is in particular upper-good, and every
upper-good routing costs at least $60$.  If $x=\sum_r\lambda_r y^{(r)}$ were
a convex combination of two-sided-good unsplittable flows, then
\[
 58=c^\top x=\sum_r\lambda_r\,c^\top y^{(r)}\ge 60,
\]
a contradiction.  The same separator applies verbatim to the cost-enhanced
two-sided statement.  The costless existence statement survives: the instance
admits two-sided-good routings, so no contradiction arises there.
\end{proof}

\begin{remark}
The underlying undirected graph is a subdivision of $K_4$, the minimal
forbidden topology for series-parallel graphs, where the convex-decomposition
statement is a theorem \cite{AlmoghrabiSkutellaWarode2026}.  The boundary is
therefore sharp: series-parallel digraphs admit the decomposition, while a
single $K_4$ subdivision already destroys it.
\end{remark}

\section{A parametric family and the planar $9/8$ bound}\label{sec:family}

Normalize $D=1$ and retain the topology of \Cref{tab:arcs}.  Let
\[
 d_1=d_3=1,\qquad d_2=b\in(0,1],
\]
let the cheap-path probabilities in the fractional path decomposition be
$p_1=p_3=r$ and $p_2=q$, and give each expensive path total integral cost
one: $c_{st_1}=1$, $c_{st_2}=1/b$, $c_{ut_3}=1$, all other costs zero.  The
fractional cost is then $3-(2r+q)$.

\begin{proposition}\label{prop:parametric}
These parameters define a strict counterexample to the additive-$D$ cost
conjecture whenever
\[
 2r+q>1,\qquad b(1-q)>r,\qquad 2r+bq<1.
\]
More generally, every no-more-expensive routing has upper deviation at least
\[
 \min\{\,1+b(1-q)-r,\ 2-2r-bq\,\}.
\]
\end{proposition}
\begin{proof}
Since $2r+q>1$, the fractional cost is strictly below $2$.  Integral routing
costs are integers, so a no-more-expensive routing uses at most one expensive
path, hence at least two cheap paths.  The pair $(Z_1,Z_2)$ exceeds the
fractional load on $su$ by $1+b(1-q)-r$; the pair $(Z_1,Z_3)$ exceeds the
load on $uv$ by $2-2r-bq$; the pair $(Z_2,Z_3)$ exceeds the load on $vw$ by
the same amount as $(Z_1,Z_2)$.  The last two displayed inequalities make
both expressions strictly larger than $D=1$.
\end{proof}

Define $\alpha_{\mathrm{planar}}$ as the infimum of constants $\alpha$ such
that every planar acyclic instance admits an unsplittable $y$ with
$c^\top y\le c^\top x$ and $y_a\le x_a+\alpha D$ for all arcs.

\begin{theorem}\label{thm:nine-eighths}
\[
 \frac98\;\le\;\alpha_{\mathrm{planar}}\;\le\;2 .
\]
Moreover, $9/8$ is the supremum of the lower bounds generated by the
symmetric family of \Cref{prop:parametric}.
\end{theorem}
\begin{proof}
The upper bound is the planar cost theorem of Traub, Vargas Koch, and
Zenklusen \cite{TraubVargasKochZenklusen2024}.  For the lower bound take
\[
 b=\frac34,\qquad r=\frac14,\qquad q=\frac12+\eps .
\]
The fractional cost is $2-\eps$, so every no-more-expensive routing uses two
cheap paths, and each of the three cheap pairs exceeds its corresponding
fractional arc load by $\tfrac98-\tfrac34\eps$.  Letting $\eps\downarrow0$
proves $\alpha_{\mathrm{planar}}\ge9/8$.

For optimality within the symmetric family, decrease $q$ to the boundary
$q=1-2r$; this can only increase both expressions in
\Cref{prop:parametric}.  If $b\ge1/2$, the two affine expressions balance at
$r=1-b$, where their common value is $b(3-2b)$, maximized at $b=3/4$ with
value $9/8$.  If $b<1/2$, neither expression exceeds $1$.  Strict cost
separation requires $q>1-2r$, so $9/8$ is approached but not attained.
\end{proof}

\section{The triangle-facet mechanism}

Let $z_i=1$ indicate that terminal $i$ uses its cheap path $Z_i$.
Capacity-good integral routings satisfy the triangle stable-set facet
$z_1+z_2+z_3\le1$, while the fractional cheap-path marginals of Rybin's
instance are $(\tfrac13,\tfrac25,\tfrac13)$ with sum $\tfrac{16}{15}>1$.  The
cost vector is a nonnegative separator for this violated facet, which is why
the same instance simultaneously kills the cost conjecture and the convex
decomposition.  As a design principle: identify binary local decisions,
characterize the capacity-feasible integral patterns as stable sets of a
conflict graph, place the fractional marginals outside the stable-set
polytope, and read the cost vector off the violated facet.

\section{Verification and disclosure}

The accompanying repository contains Rybin's instance in machine-readable
form, an exhaustive exact verifier over all eight routings
(\texttt{verify\_dgg.py}), and an enumerator for the parametric family
(\texttt{search\_dgg\_family.py}) confirming that $9/8$ is the supremum over
small-denominator parameters.  A generative AI system (GPT-5.6 Pro) assisted
with exploration, checking, code drafting, and exposition during the research
session in which these observations were made; it is not an author.  The
author assumes responsibility for verification, attribution, and the final
text.  The base counterexample is Rybin's; this note claims only the
corollaries, the parametric analysis, and the exposition.

\bibliographystyle{plainnat}
\bibliography{references}

\end{document}
